\chapter{Introduzione}

Il presente capitolo fornisce una panoramica del progetto di tirocinio, durato trecento ore, offrendo una spiegazione sull'ambito del progetto e indicando le tecnologie impiegate per la sua realizzazione.

\section{Introduzione al progetto}

Il progetto di tirocinio mira allo sviluppo e l'integrazione di un prototipo di applicazione per la visualizzazione di mappe geospaziali,
avvalendosi degli standard dell'\textit{Open Geospatial Consortium} (\textit{OGC}) e delle più moderne tecnologie web.
\\L'iniziativa nasce dalla crescente importanza dei dati geospaziali nelle applicazioni web, e della necessità di fornire strumenti interattivi, e intuitivi, per la loro visualizzazione e manipolazione.
\\Il lavoro svolto, della durata di circa trecento ore, si concentra sia sull'implementazione di un'applicazione front-end per la visualizzazione di mappe geospaziali, includendo un'interfaccia utente
interattiva, sia sullo sviluppo di un back-end per la gestione e il fornimento dei dati delle mappe.
\\Lo sviluppo si è concentrato su un'applicazione già esistente, denominata \textit{Inspicio}, creata dall'azienda Geckosoft in collaborazione con le facoltà di Informatica e di Ingegneria Informatica dell'Università di Pisa.
\\Tale applicazione nasce con l'obiettivo di creare una piattaforma completa per il censimento digitale e la valutazione del rischio dei ponti in tutta Italia
\footnote[1]{Con censimento digitale di un ponte si intende la raccolta, organizzazione e manutenzione di tutti i dati relativi a tale ponte, con lo scopo di tenerne sotto controllo lo stato in maniera continuativa.}.
La piattaforma offre agli utenti la possibilità di raccogliere dati dettagliati sui ponti, compresi parametri come posizione geografica, dimensioni, materiali di costruzione, età e condizioni attuali.
Attraverso una mappa interattiva dell'Italia, è possibile esplorare dettagliatamente tali dati.
Talvolta però ci sono delle informazioni, relative a un ponte, che non dipendono tanto dal ponte stesso quanto dalla zona ove è eretto. Tra queste troviamo: possibile attività sismiche, pericolo alluvioni, esondazione di fiumi, frane, etc.. Per tale motivo, vengono implementate delle mappe che descrivono questi fenomeni, e che possono essere sovrapposte per ottenere un quadro più generale del ponte considerato.
% Per la valutazione dei rischi strutturali, è stato impiegato un modello di intelligenza aritficiale che, sulla base di immagini raccolte in loco, fa una valutazione dell'integrità del viadotto.

\section{Base di partenza}
All'inizio del tirocinio, l'architettura del software e lo stack di sviluppo per la piattaforma Inspicio erano già stati selezionati. Tuttavia, la componente che risultava totalmente assente era quella relativa alle mappe, senza né front-end né back-end definiti. Per questo motivo, il candidato, assieme al tutore aziendale, si sono occupati di valutare e scegliere le tecnologie che meglio avrebbero consentito la gestione di mappe geospaziali. Dopodiché il candidato è stato incaricato dello sviluppo di questo aspetto. La precedente mancanza delle mappe si riconduce ad altre priorità degli sviluppatori già presenti in azienda.

\section{Contesto di Sviluppo}
La piattaforma, che prende l'acronimo di \textit{BMS} (\textit{Bridge Management System}) è stata sviluppata come un'applicazione Web moderna, seguendo un'architettura a microservizi che favorisce la modularità e la scalabilità del sistema. Questa scelta architetturale è stata preferita rispetto alle tradizionali applicazioni monolitiche, in quanto ha permesso di separare i vari servizi, semplificandone la gestione e consentendo una maggiore flessibilità nell'implementazione di nuove funzionalità.
Ciò ha portato all'adozione di un ambiente di sviluppo basato su un approccio containerizzato utilizzando \textit{Docker}. L'utilizzo di container Docker ha semplificato notevolmente la gestione delle dipendenze e delle configurazioni dei servizi, garantendo al contempo una maggiore portabilità dell'applicazione tra diversi ambienti di sviluppo e produzione.
\\Per quanto concerne le tecnologie impiegate nello sviluppo dell'interfaccia grafica della piattaforma, è stato selezionato \textit{Angular}, un framework JavaScript moderno e robusto, caratterizzato da una struttura modulare che offre una vasta gamma di librerie. Queste caratteristiche lo rendono particolarmente adatto per la creazione di interfacce utente dinamiche e complesse, come nel caso della piattaforma Inspicio.
\\Per il back-end, è stata selezionata la piattaforma .NET, un framework di sviluppo completo e versatile supportato da Microsoft. .NET offre una vasta gamma di strumenti e librerie che semplificano lo sviluppo, la manutenzione e la scalabilità delle applicazioni.
\\L'adozione di Angular e .NET come principali framework di sviluppo hanno consentito agli sviluppatori di seguire un pattern unificato attraverso una varietà di progetti. Questa scelta ha ridotto la confusione e garantito una maggiore coerenza nelle implementazioni, assicurando che tutti i software sviluppati seguissero gli stessi standard e convenzioni. Si tratta di un aspetto cruciale in contesti aziendali con molti progetti in corso, facilitando la gestione e la manutenzione dell'intero insieme di software.
\\Infine, nella comunicazione tra front-end e back-end della piattaforma, è stato adottato il paradigma delle \textit{Rest API}, che rappresenta uno standard per lo scambio di dati tra client e server su protocollo HTTP. Più specificamente, le API sono esposte tramite Swagger UI, consentendo al front-end di interagire con esse in modo intuitivo e documentato. Il protocollo utilizzato per definire le specifiche delle API rispecchia lo standard di OpenAPI 3 (OAS3), garantendo una chiara e precisa descrizione delle operazioni e dei dati scambiati tra front-end e back-end.

\section{Struttura della Relazione}

\begin{description}
    \item[{\hyperref[cap:chapter2]{Il secondo capitolo}}] descrive le tecnologie studiate e utilizzate per lo sviluppo del progetto.
    \item[{\hyperref[cap:chapter2]{Il terzo capitolo}}] approfondisce il percorso di tirocinio formativo.
          % \item[{\hyperref[cap:chapter2]{Il quarto capitolo}}] approfondisce il progetto da un punto di vista tecnico, descrivendo i requisiti e i casi d'uso del progetto.
          % \item[{\hyperref[cap:chapter2]{Il quinto capitolo}}] approfondisce le tecnologie utilizzate per la realizzazione del progetto, descrivendo le scelte progettuali e le soluzioni implementate, a livello architetturale e di codifica.
          % \item[{\hyperref[cap:chapter2]{Il sesto capitolo}}] approfondisce le attività di verifica e validazione del progetto, descrivendo le tipologie di test effettuate e i risultati ottenuti.
          % \item[{\hyperref[cap:chapter2]{Nel settimo capitolo}}] descrive lo stato degli obiettivi raggiunti e la valutazione personale del percorso, con un'analisi retrospettiva del progetto e del suo utilizzo rispetto a sviluppi futuri sulle stesse tematiche.
\end{description}